Hlavním cílem této diplomové práce bylo navrhnout autonomního agenta, který je schopen fungovat v~dynamickém herním prostředí s~využitím metod strojového učení a~adaptovat své chování, včetně volby výbavy, s~cílem efektivně čelit protihráči.

V~úvodních kapitolách byly představeny existující přístupy k~umělé inteligenci ve hrách a~nejčastěji používané algoritmy a~techniky. Mezi analyzovanými systémy se nacházely mimo jiné AlphaStar, OpenAI Five a~Voyager; z~metod pak zejména Proximal Policy Optimization (PPO), Deep Q-Network (DQN), Self-Play a~Population-Based Training.

Následně byly popsány technologie použité při realizaci samotné hry a~agenta, především herní engine Unreal Engine a~plugin NevarokML. Součástí této části byl rovněž návrh herní logiky a~architektury agenta.

Závěrečná část práce se zaměřila na implementaci jednotlivých herních modulů, jako jsou hráčská postava, zbraně, systém kol a~mechanismy získávání výbavy, a~především na návrh, trénování a~vyhodnocení jednotlivých verzí agenta. Výsledný agent je schopen pohybu, rotace, útoku i~změny výbavy a~jeho chování je řízeno neuronovou sítí trénovanou metodami posilovaného učení.

První verze agenta byla záměrně jednoduchá a~sloužila především k~ověření práce s~pluginem NevarokML a~celým trénovacím frameworkem. Druhá verze již obsahovala kompletní bojovou logiku, avšak vykazovala nestabilní a~místy nepřirozené chování, zejména v~oblasti udržování vzdálenosti od protivníka. Finální verze agenta tyto nedostatky ve velké míře odstranila a~nejlépe prokázala schopnost adaptivního chování v~dynamickém prostředí. Agent je schopen aktivně pronásledovat protivníka, útočit a~po jednotlivých kolech upravovat svou výbavu na základě předchozích střetů.

\newpage

Navzdory dosaženým výsledkům se ukázalo, že skutečná adaptace agenta v~reálném čase během samotného hraní je výrazně omezená. Tento stav je způsoben především omezeními pluginu NevarokML. V~jeho současné podobě je například možné provozovat pouze jednoho trenéra na jedno dynamicky vytvářené herní prostředí, zatímco u~statických trénovacích map tento problém nenastává, jelikož je možné přiřadit jednomu trenérovi více prostředí. Dále není podporována změna odměnové funkce, akčního prostoru ani vstupních pozorování v~průběhu tréninku při zachování již naučeného modelu, tedy přístup známý jako Curriculum Learning, založený na postupném zvyšování obtížnosti úlohy. Z~tohoto důvodu je nutné trénovat jeden komplexní model od počátku na plně složité úloze, což výrazně omezuje možnosti dosažení optimálního chování agenta.